\documentclass[11pt,a4paper]{article}
\usepackage[margin=1in]{geometry}
\usepackage{booktabs}
\usepackage{hyperref}
\usepackage{xurl}

\title{Step5a-HGAT Paper Prep: Algorithm and Experiment Snapshot}
\author{}
\date{\today}

\begin{document}
\maketitle

\section{Goal}
This draft consolidates algorithm and ablation evidence before full paper writing.
Our focus is to verify whether the proposed topology-conditioned architecture improves OOD generalization under table-group split.

\section{Algorithm Core: TCDR}
We use a shared regression backbone and add topology-conditioned dual residual experts:
\begin{itemize}
  \item shared base head captures cross-domain common trends;
  \item topology embedding conditions two domain-specific residual heads (src/tgt);
  \item final prediction = base prediction + domain residual.
\end{itemize}

This decouples universal delay behavior from topology/domain-specific bias.

\section{Dataset and Protocol}
Primary dataset for paper tables:
\begin{center}
\path{output/dataset_table_group_1166_with_test.pkl}
\end{center}

Target split statistics:
Train groups = 25 (1225 samples), Val groups = 25 (1225 samples), Test groups = 104 (5094 samples).

Unified training setup in this round:
batch size 2048, learning rate $2\times10^{-3}$, frozen HGAT encoder, parasitic append\_split features, 3 seeds (9294/9295/9296).

\section{Main Results (3-seed Mean$\pm$Std)}
\begin{table}[htbp]
\centering
\small
\begin{tabular}{lccc}
\toprule
Method & Val R2 (mean+/-std) & Test R2 (mean+/-std) & Delta Test vs Baseline \\
\midrule
Baseline (shared linear) & 0.8554+/-0.0024 & 0.8496+/-0.0045 & +0.0000 \\
TCDR (16,64) & 0.8859+/-0.0082 & 0.8888+/-0.0056 & +0.0393 \\
TCDR (16,96) & 0.8843+/-0.0097 & 0.8860+/-0.0061 & +0.0364 \\
TCDR (32,64) & \textbf{0.8931+/-0.0098} & \textbf{0.8951+/-0.0099} & \textbf{+0.0455} \\
\bottomrule
\end{tabular}
\caption{Core architecture comparison under table-group split.}
\end{table}

\section{Component Ablation around TCDR(16,64)}
\begin{table}[htbp]
\centering
\small
\begin{tabular}{lcc}
\toprule
Method & Val R2 (mean+/-std) & Test R2 (mean+/-std) \\
\midrule
TCDR(16,64) & 0.8859+/-0.0082 & 0.8888+/-0.0056 \\
+ dual\_readout(mean) & 0.8859+/-0.0082 & 0.8888+/-0.0056 \\
+ calib\_mlp & \textbf{0.8951+/-0.0104} & 0.8930+/-0.0070 \\
+ parasitic\_replace & 0.8716+/-0.0116 & 0.8842+/-0.0108 \\
\bottomrule
\end{tabular}
\caption{Component-level ablation results.}
\end{table}

\section{Takeaways}
\begin{itemize}
  \item TCDR delivers a clear and consistent gain over baseline on OOD split.
  \item TCDR(32,64) is best in mean test R2 for this round.
  \item TCDR(16,64) remains a strong and simpler stable candidate.
  \item dual\_readout gives negligible gain in current setting; parasitic\_replace is less robust.
\end{itemize}

\section{Artifact Files}
\begin{itemize}
  \item Per-run logs and metrics: \path{copilot_train_logs/paper_main_ablation_runs.csv}
  \item Aggregated summary: \path{copilot_train_logs/paper_main_ablation_summary.csv}
  \item Paper table markdown: \path{copilot_train_logs/paper_main_ablation_table.md}
\end{itemize}

\end{document}
